\documentclass{article}
\usepackage{arxiv}

\usepackage[utf8]{inputenc}
\usepackage[T1]{fontenc}
\usepackage{hyperref}
\usepackage{url}
\usepackage{booktabs}
\usepackage{amsmath, amssymb}
\usepackage{bm}
\usepackage{natbib}
\usepackage{graphicx}

% The emblem lives in the repository-wide assets directory, two levels up.
\graphicspath{{../../assets/}{../../assets/logos/}{../../assets/biota/}{../../assets/diagrams/}}

\title{Support Modules and the False Alarm Budget}

% arxiv.sty labels every paper "A Preprint" in the running head and under the
% title. This is an internal technical report, not a preprint of a submission.
\renewcommand{\headeright}{Technical Report}
\renewcommand{\undertitle}{Technical Report}

% arxiv.sty's title block sets a 24pt strut above the author names and 0.4in
% below them, which leaves the byline stranded between "Technical Report" and
% the date. This is the same layout with those two gaps halved. Overridden here
% rather than edited in arxiv.sty, which is a vendored file.
\makeatletter
\renewcommand{\@maketitle}{%
  \vbox{%
    \hsize\textwidth
    \linewidth\hsize
    \vskip 0.1in
    \@toptitlebar
    \centering
    {\LARGE\sc \@title\par}
    \@bottomtitlebar
    \textsc{\undertitle}\\
    \vskip 0.1in
    \def\And{%
      \end{tabular}\hfil\linebreak[0]\hfil%
      \begin{tabular}[t]{c}\bf\rule{\z@}{12\p@}\ignorespaces%
    }
    \def\AND{%
      \end{tabular}\hfil\linebreak[4]\hfil%
      \begin{tabular}[t]{c}\bf\rule{\z@}{12\p@}\ignorespaces%
    }
    \begin{tabular}[t]{c}\bf\rule{\z@}{12\p@}\@author\end{tabular}%
  \vskip 0.2in \@minus 0.1in \center{\@date}   \vskip 0.2in
  }
}
\makeatother

\author{
  Daniel Boupha Christensen \\
  Environmental, Health, and Safety, University of California, Berkeley \\
  The Green Initiative Fund \\
}

\begin{document}

% Masthead. Licence at the left, group emblem at the right. This is an internal
% report, so the block carries the licence and the project mark only, and
% deliberately no journal, volume, page range or DOI furniture.
\noindent
\begin{minipage}[c]{0.74\textwidth}
  \footnotesize
  Copyright 2024--2026 Strawberry Creek Monitoring Group contributors.\\
  Licensed under the Apache License, Version 2.0.\\
  \url{https://www.apache.org/licenses/LICENSE-2.0}
\end{minipage}%
\hfill
\begin{minipage}[c]{0.16\textwidth}
  \raggedleft
  \includegraphics[width=\linewidth]{SCMGemblem.png}
\end{minipage}

\maketitle

\begin{abstract}
Support modules are optional layers that attach to a detection run. They are
not models and have no weights of their own, and one module attaches to any
model that declares the input contract it needs. The design problem they
address is narrower than it may appear. Attaching a second test to a detector
calibrated at a false alarm rate $q$ produces a system that alerts at roughly
$2q$, while every component still reports $q$ and nothing in the arithmetic
signals the change. Here, we set out (1) the taxonomy that makes the false
alarm budget divisible, (2) the algebra that divides it, (3) the audit record
that states which module caused which alert, and (4) the mathematics of Trial
Bed, the first module we implemented. Trial Bed is an explainer. It names the
disturbance a firing resembles by matching the direction each channel moved
against a signature table, and it declines to answer when the channels a
signature needs are not measured. Since the first version of this report, we
have built Newtwork Run as a second explainer, which reads the order in which
sites fired along the flow graph, and Settling Pool as a screen, which flags
readings that are not measurements. All three can be attached to a model in
the StrawberryWatch dashboard.
\end{abstract}

\section{Introduction}
\label{sec:introduction}

We calibrate a detector once, against a fault-free distribution, at a false
alarm rate chosen in advance. That calibration is what keeps the alert stream
small enough for an operator to read, and it has a specific weakness: it
describes one test. Adding a second test beside it removes the guarantee
without any warning, because each test is individually correct.

We built the support layer around this failure. If a primary detector fires
with probability $q$ under the null and a second detector fires with
probability $q$ under the null, the system fires with probability up to $2q$.
Both components report $q$ when asked. No exception is raised and no number
looks wrong, yet the alarm rate has doubled.

Thus, a false alarm budget is a quantity to be divided, and it is never handed
out afresh to each module. The rest of this document follows from this
principle.

\section{Module Kinds}
\label{sec:kinds}

Not every module that attaches to a run costs budget, and treating the modules
that do in the same way as those that do not would make the budget
indivisible. The layer therefore has four kinds, and each module is exactly one
of them. Taken in order, the kinds follow the order of the pipeline.

\begin{table}[t]
  \centering
  \caption{The four kinds. Only one of them spends false alarm budget.}
  \label{tab:kinds}
  \begin{tabular}{lllll}
    \toprule
    Kind & Runs & Method & Changes & Budget \\
    \midrule
    screen    & before the model & \texttt{admit}       & the model's input & none \\
    modulator & alongside        & \texttt{multipliers} & the threshold     & none \\
    detector  & after            & \texttt{score}, \texttt{null} & the verdict & yes \\
    explainer & after the decision & \texttt{explain}   & nothing           & none \\
    \bottomrule
  \end{tabular}
\end{table}

A \textbf{modulator} moves the alerting threshold and never touches the score.
This direction has a statistical basis. The threshold was fitted against the
fault-free distribution of the score, so a score scaled by any factor is no
longer a draw from the distribution the fit was made on, and every probability
computed under it loses its meaning. Multiplying the threshold instead leaves
the score, its probability and the calibration unchanged, and it makes the
adjustment a stated alerting decision that can be read from the record.

A \textbf{detector} adds a score, and therefore a test and a second tail to pay
for. It is the only kind that spends budget, and separating it from the
modulator is what makes the budget divisible. With one class and a flag,
nothing would prevent three modules from each running at $q$ while the real
rate rose to $3q$.

A \textbf{screen} runs before the model and decides which readings are fit to
use. It adds no test, so it is not a detector and takes no budget, and it moves
no threshold, so it is not a modulator. It changes the model's input, which has
a consequence that neither of the other kinds has. The null and the threshold
were fitted on unscreened windows, so screening moves the distribution the
score is drawn from, and the existing calibration no longer describes it. A
screen therefore cannot be switched on without refitting the primary.

An \textbf{explainer} runs after the decision has been made and reads it. It
adds no test, moves no threshold and changes nothing the decision was made
from. It is a separate kind rather than a detector because charging it budget
would cost the primary sensitivity in exchange for a module that can never
fire.

\section{The False Alarm Budget}
\label{sec:budget}

\subsection{Division of $q$}

Let $q$ be the total false alarm rate for the whole run, and let there be $n$
attached detectors with weights $w_{1}, \dots, w_{n}$. The split is
\begin{equation}
  q_{\mathrm{primary}} = q \left( 1 - \sum_{i=1}^{n} w_{i} \right),
  \qquad
  q_{i} = q\, w_{i} .
\end{equation}

This is a union bound. The probability that at least one of the tests fires
under the null is at most the sum of the individual rates, so allocating shares
that sum to $q$ keeps the combined rate at or below $q$ rather than at $q$. It
is the same correction used whenever a family of tests has to hold a joint
error rate~\citep{dunn1961multiple}, applied here to a family that grows
whenever an operator adds a module on the command line.

The default is an even split. With $n$ detectors there are $n+1$ shares, so
each weight is $1/(n+1)$ and the weights sum to $n/(n+1)$, which is below one
for every $n$. We make the even split the default because weighting one
detector above another is a claim about which test is more likely to fire, and
nothing has measured that.

\subsection{The requirement that weights sum below one}

We require the weights to sum strictly below one, and both boundary cases raise
an error rather than being clamped. At a sum of exactly one, the primary is
left a budget of zero, which is a threshold at infinity and a primary that can
never fire. Above one, the primary's budget is negative, which is not a rate.

Clamping either case would produce a run that appears to have worked, reports a
plausible split, and has switched the primary detector off without notice.
Raising is the only safe behaviour, because otherwise the failure is invisible
in the way this layer exists to prevent.

Modulators, screens and explainers do not enter this arithmetic. They add no
test, so $n$ counts detectors only, and a stack of modulators alone leaves the
primary the whole budget. At present no detector module is registered, so
every stack leaves the primary the whole budget.

\subsection{Firing on the probability}

Each detector fires when its right-tail probability falls below its own share,
$p_{i} < q_{i}$, and never on the raw score. Raw scores from different
detectors are not comparable, for the same reason that a multi-channel detector
converts each channel before combining them. Comparing on the probability keeps
the meaning of the budget the same for every module, whatever units its score
is in.

We fit a detector's null at calibration time and load it from an artifact; the
support layer never fits it. A detector that fits its null on the window it is
judging does not have a null; it has a description of the data it just saw.
Where a detector's null is a fitted tail, it uses the same peaks-over-threshold
construction as the primary~\citep{siffer2017spot,pickands1975gpd}, and where
several probabilities have to become one number, it uses the same combination
rule~\citep{fisher1925methods}.

\section{Composition}
\label{sec:composition}

\subsection{The stack}

The modules attached to one run form a stack. A loader builds the stack, so
that name resolution, contract validation, collision checks against a model's
built-in behaviour and the budget split all happen before anything runs. Thus,
a stack that exists has already been checked, and a module that cannot attach
never appears partway through a detection.

Screens intersect. A reading is admitted only if every screen admits it, and an
empty screen list admits everything, which is the identity.

Modulator multipliers compose by product,
\begin{equation}
  M_{t,j} = \prod_{m} M^{(m)}_{t,j},
  \qquad M^{(m)}_{t,j} \ge 1 ,
\end{equation}
and every multiplier must be finite and at least one. This layer raises the
alerting bar and never lowers it, so a factor below one is refused rather than
applied. A timestep with no reading behind it contributes exactly one, so a
modulator without information is inert rather than suppressing alerts.

\subsection{Ordering}

We sort modules by name before composition, so the order in which an operator
typed them cannot affect anything downstream. We sort rather than rely on
commutativity for a specific reason. Multiplication commutes, so $ab = ba$
holds, but floating point association does not: $(ab)c$ and $(ac)b$ need not
be bit-identical, and three modulators are enough to reach that case. An empty
stack is the identity with no special case, so zero modules follow the same
code path as three.

\section{The Audit Record}
\label{sec:record}

The record answers two questions an operator asks: why an alert fired, and
what would change if a given module were turned off.

\subsection{Three verdicts}

Every node at every timestep resolves to one of three verdicts: it fired, it
was clear, or it was not scorable. The third verdict exists because a score
that could not be computed is a different answer from a score that did not
exceed the threshold, and merging the two would let an unscorable node reach an
operator as an all-clear. A non-finite score never fires, and where nothing
was scorable the record says so.

\subsection{Counterfactuals}

For each attached module, the record carries the verdict that would have been
reached with that module switched off, computed without re-running anything. A
modulator drops out by dividing its multiplier back out of the threshold, and a
detector drops out by removing its vote. An explainer's entry is written too,
and it always equals the actual verdict, so the record itself shows that the
explainer cost nothing.

The record extends the primary's own format rather than replacing it. Every
field the primary wrote keeps its meaning, and the support fields sit
alongside them. Two formats would require two readers, and an operator who has
to act on an alert at two in the morning should not have to work out which
format they are reading.

\section{Trial Bed}
\label{sec:trialbed}

Trial Bed is the first module we implemented, and it is an explainer. Given a
firing, it names the disturbance the event resembles, the evidence behind that
name, and how much confidence the reporting channels support.

\subsection{Directions rather than magnitudes}

For each parameter, the module compares the readings at the timesteps that
fired against the readings at the timesteps that did not, and reduces the
comparison to a direction. Writing $\bar{b}$ and $s_{b}$ for the baseline mean
and standard deviation and $\bar{e}$ for the event mean,
\begin{equation}
  \text{direction} =
  \begin{cases}
    \textsc{up}   & (\bar{e} - \bar{b}) / s_{b} \ge \kappa, \\
    \textsc{down} & (\bar{e} - \bar{b}) / s_{b} \le -\kappa, \\
    \textsc{flat} & \text{otherwise},
  \end{cases}
  \qquad \kappa = 1 .
\end{equation}

Using the baseline's own spread as the unit lets one rule serve parameters on
very different scales. A move of one unit in temperature and a move of one
unit in conductivity are not comparable in absolute terms, but both are
comparable in units of how much that parameter normally varies at that node.
Where the baseline has no spread at all, the comparison falls back to a
relative change with a five percent band, and where there is too little data,
the direction is left undefined rather than guessed.

\subsection{The signature table}

Table~\ref{tab:signatures} shows the deployed table. Each cell is the direction
a disturbance is expected to push a parameter, and \textsc{indet} marks a
parameter that does not discriminate for that disturbance and is therefore
skipped, with neither reward nor penalty.

\begin{table}[t]
  \centering
  \caption{The disturbance signature table as deployed. \textsc{indet} means the
    parameter does not discriminate for that disturbance. Entries marked $\ast$
    are expected to move strongly; the current matcher does not yet weight
    them differently.}
  \label{tab:signatures}
  \begin{tabular}{llllll}
    \toprule
    & temperature & dissolved oxygen & pH & conductivity & floating cond. \\
    \midrule
    rain       & \textsc{indet} & up   & down & down & down \\
    potable water   & down           & down & up   & down & down \\
    oil        & up             & down$^{\ast}$ & down & down & down$^{\ast}$ \\
    sewage     & up             & down$^{\ast}$ & \textsc{indet} & up & up \\
    fertilizer & flat           & down & \textsc{indet} & up & up \\
    \bottomrule
  \end{tabular}
\end{table}

Each cell also carries the mechanism behind its direction, for example that
oil is a poor conductor for the oil row's conductivity cell, and every
per-parameter comparison in an account quotes it. A row is scored by the
fraction of its non-\textsc{indet} parameters whose observed direction agrees,
\begin{equation}
  \mathrm{score}(P) = \frac{|\{v \in R(P) : d_{v} = \sigma_{P,v}\}|}{|R(P)|},
\end{equation}
where $R(P)$ is the set of parameters for which row $P$ asserts a direction. A
row scoring below one half is a poor fit and is reported as a possible new type
rather than as a name. A tie at the top is reported the same way, since if the
reporting channels did not separate two candidates, naming one of them would
invent the separation.

The potable water row is also checked against numbers rather than directions
alone. A move down is weaker evidence than a mean of about 13 degrees Celsius
when the creek already runs near 13 degrees, so the module compares the event
means with Berkeley potable water, about 13.0 degrees Celsius and pH 9.4, and
reports a match within 1.5 degrees and 0.5 pH units. This check is reported
beside the row's score and does not change it.

\subsection{Declining to answer}

The most important property of the module is its behaviour when the channels a
row needs are not measured. A row is evaluated only if every parameter it
asserts a direction for has an observed direction behind it. Otherwise, the row
is reported as unevaluable, together with the names of the absent parameters
and, from the inventory, the reason each one is absent: a probe that was never
installed at the site, a probe that is installed but reporting nothing, or a
probe with no readings in the window. The first two look identical in the data
and are different problems for whoever receives the alert.

The alternative, scoring a row on whichever of its parameters happen to be
present, produces a number that looks like a full match and is not one. A
partial signature scored as if it were complete would present a guess to an
operator as a diagnosis.

This concern applies directly to this network. Depth is not in the table at
all, every row in Table~\ref{tab:signatures} requires both dissolved oxygen
and floating conductivity, and four of the five rows also require pH. Of the
11 sites in the inventory, 9 measure only conductivity, depth and temperature,
and no row is evaluable there. Footbridge and south\_fork\_3 also carry
dissolved oxygen and floating conductivity probes, so the sewage and
fertilizer rows, which do not depend on pH, can be evaluated at those two
sites. No site has a pH probe, so the rain, potable water and oil rows are
evaluable nowhere. The four sites DuskCrayfish scores are among the nine, so
when Trial Bed is attached to DuskCrayfish its output at every node is that it
cannot evaluate, together with the list of channels that would let it do so.

\subsection{The output}

An account states the likely cause, the evidence and the confidence, in a form
an operator can act on rather than as a bare label. Where a diagnosis is
possible, it reads, for example, \emph{consistent with sewage: temperature up,
dissolved oxygen down, pH never installed here, conductivity up, floating
conductivity up}. Where it is not, the same sentence structure carries the refusal and the
missing channels. We report confidence from how many of the discriminating
channels were populated rather than from the score, because a perfect match on
two channels is not strong evidence.

\section{Newtwork Run and Settling Pool}
\label{sec:pending}

When the first version of this report was written, both modules were
registered placeholders. Newtwork Run was planned as a detector, which would
score agreement along the flow edges against a fitted null and spend a share
of the budget, and Settling Pool was planned as a screen. Both are now built.

\subsection{Newtwork Run}

We built \textbf{Newtwork Run} as an explainer rather than a detector. A real
discharge appears at a downstream site after it appears upstream, while a
catchment-wide cause such as rain moves every site at once, and the order in
which sites first fired already separates the two without a fitted null or a
measured travel time. For each node that fired, the module compares its first
firing with those of its upstream and downstream neighbours in the inventory's
flow graph and returns one of four readings. (1) A site upstream fired
earlier, or a site downstream fired later, which is consistent with transport.
(2) Other sites fired at the same timestep with no upstream lead, which is what
a catchment-wide cause looks like. (3) No neighbour fired before or after it along
the flow graph and no site fired at the same timestep, so nothing suggests the
event moved through the creek. (4) Codornices, which is not connected to
Strawberry Creek by flow, fired at the same timestep, so nothing carried the
event between them. The table of travel times per edge is left empty, because
nobody has measured one, and every account states that its confidence rests
on ordering alone. As an explainer, the module spends no budget and cannot
change what fired.

\subsection{Settling Pool}

\textbf{Settling Pool} is a screen. It runs eight tests on every site and
variable, with thresholds per site from the inventory: (1) gross range,
(2) flat line, (3) attenuated signal, (4) spike, (5) rate of change,
(6) sentinel values such as $-9999$, (7) a duplicate feed, where a site reports
another site's numbers, and (8) staleness. Each reading receives the worst flag
any test gave it, on five levels numbered as in the QARTOD quality control
manuals: 1 good, 2 not evaluated, 3 suspect, 4 fail and 9 missing. The
inventory's record of each sensor outranks the tests where it says a reading
was not a measurement, so a probe that was not installed is not evaluated
rather than tested, and a site taken out of service is withheld whatever its
flags say. Settling Pool flags readings and does not correct them. What a
model does with each level is that model's flag policy, and the default policy
refuses fail and missing.

In the dashboard, an operator can attach Settling Pool to a model. Readings
flagged suspect or fail are then blanked before the model reads them, and the
model's account reports how many readings were held back at each site. As
Section~\ref{sec:kinds} sets out, a screen moves the distribution the primary
was calibrated against, so switching it on still calls for recalibrating the
primary on screened windows.

\section{Limitations and Future Work}
\label{sec:limitations}

The budget split is a union bound, which is conservative. The real combined
rate falls below $q$ rather than at $q$, and how far below depends on how
correlated the tests are. Nothing has yet measured that correlation for this
system, so the amount of sensitivity given up is currently unknown.

The even default split reflects the absence of a measurement rather than a
finding. It is the appropriate default only until something measures which
test is more likely to fire. With no detector module registered, the split is
currently unused.

No model has yet been recalibrated on screened windows. Attaching Settling Pool
in the dashboard therefore scores screened readings against thresholds fitted
on unscreened ones. In addition, the dashboard blanks readings flagged suspect
or fail, while the default flag policy that models declare refuses fail and
missing. The two should be brought into agreement.

Within Table~\ref{tab:signatures}, floating conductivity has the same direction
as conductivity in all five rows. It nonetheless counts as a discriminating
channel, and one discriminating channel is currently enough to unlock a named
diagnosis. It also means that conductivity's evidence is counted twice in every
row's score. Whether this is intended is a question about the table, and the
table belongs to the Water Quality team.

Every cell now states the mechanism behind its direction, but no cell cites
a measurement or a publication. The table is attributed as a whole, and the
two measurable quantities it relies on, the local potable water temperature
of 13.0 degrees Celsius and pH of 9.4, have nothing recorded behind them.
One labelled event also contradicts a row. In the March 2026 hydrant break,
water from a fire hydrant ran down Euclid Avenue into a storm drain and into
the north fork, and north\_fork\_0 saw conductivity fall and depth rise with
no meaningful change in temperature, where the potable water row expects
temperature to fall. Either street runoff warms on the way in, or the row is
wrong for water delivered through a drain. We recorded the event as a known
exception beside the table and left the row as sourced, since the question
belongs to the group. Depth is measured
at every deployed site and appears nowhere in the table, which is part of why
no row is currently evaluable.

Finally, Trial Bed reduces a window to one direction per parameter, which
discards magnitude and timing entirely. A signature table that distinguished a
large move from a small one, or an abrupt move from a gradual one, would need
both a richer table and a matcher that reads it, and neither exists yet.

\bibliographystyle{unsrtnat}
\bibliography{references}

\end{document}
